% Options for packages loaded elsewhere
\PassOptionsToPackage{unicode}{hyperref}
\PassOptionsToPackage{hyphens}{url}
\documentclass[
]{article}
\usepackage{xcolor}
\usepackage{amsmath,amssymb}
\setcounter{secnumdepth}{-\maxdimen} % remove section numbering
\usepackage{iftex}
\ifPDFTeX
  \usepackage[T1]{fontenc}
  \usepackage[utf8]{inputenc}
  \usepackage{textcomp} % provide euro and other symbols
\else % if luatex or xetex
  \usepackage{unicode-math} % this also loads fontspec
  \defaultfontfeatures{Scale=MatchLowercase}
  \defaultfontfeatures[\rmfamily]{Ligatures=TeX,Scale=1}
\fi
\usepackage{lmodern}
\ifPDFTeX\else
  % xetex/luatex font selection
\fi
% Use upquote if available, for straight quotes in verbatim environments
\IfFileExists{upquote.sty}{\usepackage{upquote}}{}
\IfFileExists{microtype.sty}{% use microtype if available
  \usepackage[]{microtype}
  \UseMicrotypeSet[protrusion]{basicmath} % disable protrusion for tt fonts
}{}
\makeatletter
\@ifundefined{KOMAClassName}{% if non-KOMA class
  \IfFileExists{parskip.sty}{%
    \usepackage{parskip}
  }{% else
    \setlength{\parindent}{0pt}
    \setlength{\parskip}{6pt plus 2pt minus 1pt}}
}{% if KOMA class
  \KOMAoptions{parskip=half}}
\makeatother
\usepackage{longtable,booktabs,array}
\newcounter{none} % for unnumbered tables
\usepackage{calc} % for calculating minipage widths
% Correct order of tables after \paragraph or \subparagraph
\usepackage{etoolbox}
\makeatletter
\patchcmd\longtable{\par}{\if@noskipsec\mbox{}\fi\par}{}{}
\makeatother
% Allow footnotes in longtable head/foot
\IfFileExists{footnotehyper.sty}{\usepackage{footnotehyper}}{\usepackage{footnote}}
\makesavenoteenv{longtable}
\usepackage{graphicx}
\makeatletter
\newsavebox\pandoc@box
\newcommand*\pandocbounded[1]{% scales image to fit in text height/width
  \sbox\pandoc@box{#1}%
  \Gscale@div\@tempa{\textheight}{\dimexpr\ht\pandoc@box+\dp\pandoc@box\relax}%
  \Gscale@div\@tempb{\linewidth}{\wd\pandoc@box}%
  \ifdim\@tempb\p@<\@tempa\p@\let\@tempa\@tempb\fi% select the smaller of both
  \ifdim\@tempa\p@<\p@\scalebox{\@tempa}{\usebox\pandoc@box}%
  \else\usebox{\pandoc@box}%
  \fi%
}
% Set default figure placement to htbp
\def\fps@figure{htbp}
\makeatother
\ifLuaTeX
  \usepackage{luacolor}
  \usepackage[soul]{lua-ul}
\else
  \usepackage{soul}
\fi
\setlength{\emergencystretch}{3em} % prevent overfull lines
\providecommand{\tightlist}{%
  \setlength{\itemsep}{0pt}\setlength{\parskip}{0pt}}
\usepackage{bookmark}
\IfFileExists{xurl.sty}{\usepackage{xurl}}{} % add URL line breaks if available
\urlstyle{same}
\hypersetup{
  hidelinks,
  pdfcreator={LaTeX via pandoc}}

\author{}
\date{}

\begin{document}

Assignment

STAT 4104

Submitted by: Mst. Fatema Tanbin Lava

ID: 12110003

Submitted to: Dr. Md. Siddikur Rahman (Associate Professor)
\textbf{\ul{\hfill\break
}}

\textbf{\ul{Problem 1:}}

\textbf{Lung Cancer Data Analysis by using python.}

\textbf{Python Code:}

{\def\LTcaptype{none} % do not increment counter
\begin{longtable}[]{@{}
  >{\raggedright\arraybackslash}p{(\linewidth - 0\tabcolsep) * \real{0.9987}}@{}}
\toprule\noalign{}
\begin{minipage}[b]{\linewidth}\raggedright
import pandas as pd

import numpy as np

import matplotlib.pyplot as plt

import seaborn as sns

import statsmodels.formula.api as smf

from scipy.stats import chi2\_contingency, shapiro

\# Load the dataset

df = pd.read\_csv(\textquotesingle lung\_disease.csv\textquotesingle)

\# Display first 5 rows

print(df.head())
\end{minipage} \\
\midrule\noalign{}
\endhead
\bottomrule\noalign{}
\endlastfoot
\end{longtable}
}

\textbf{(Basic EDA Questions)}

\textbf{i) Summarize the distribution of Age}

Python code:

{\def\LTcaptype{none} % do not increment counter
\begin{longtable}[]{@{}
  >{\raggedright\arraybackslash}p{(\linewidth - 0\tabcolsep) * \real{0.9987}}@{}}
\toprule\noalign{}
\begin{minipage}[b]{\linewidth}\raggedright
print(df{[}\textquotesingle Age\textquotesingle{]}.describe())

print(\textquotesingle Skewness =\textquotesingle,
df{[}\textquotesingle Age\textquotesingle{]}.skew())
\end{minipage} \\
\midrule\noalign{}
\endhead
\bottomrule\noalign{}
\endlastfoot
\end{longtable}
}

\textbf{Output:}

{\def\LTcaptype{none} % do not increment counter
\begin{longtable}[]{@{}
  >{\raggedright\arraybackslash}p{(\linewidth - 2\tabcolsep) * \real{0.5000}}
  >{\raggedright\arraybackslash}p{(\linewidth - 2\tabcolsep) * \real{0.5000}}@{}}
\toprule\noalign{}
\begin{minipage}[b]{\linewidth}\raggedright
count
\end{minipage} & \begin{minipage}[b]{\linewidth}\raggedright
500.000000
\end{minipage} \\
\midrule\noalign{}
\endhead
\bottomrule\noalign{}
\endlastfoot
mean & 43.904000 \\
std & 14.604841 \\
min & 20.000000 \\
25\% & 30.750000 \\
50\% & 45.000000 \\
75\% & 56.000000 \\
max & 69.000000 \\
Skewness: & 0.0121 \\
\end{longtable}
}

\textbf{Interpretation:} The distribution is nearly symmetric (not
significantly skewed), it does not perfectly follow a normal
distribution.

ii) Proportion of individuals (Smokers vs Non-smokers):

Python code:

{\def\LTcaptype{none} % do not increment counter
\begin{longtable}[]{@{}
  >{\raggedright\arraybackslash}p{(\linewidth - 0\tabcolsep) * \real{0.9987}}@{}}
\toprule\noalign{}
\begin{minipage}[b]{\linewidth}\raggedright
smoke\_prop =
df{[}\textquotesingle Smoking\textquotesingle{]}.value\_counts(normalize=True)
* 100

print(f"\textbackslash nSmoking Proportions
(\%):\textbackslash n\{smoke\_prop\}")
\end{minipage} \\
\midrule\noalign{}
\endhead
\bottomrule\noalign{}
\endlastfoot
\end{longtable}
}

\textbf{Output:}

{\def\LTcaptype{none} % do not increment counter
\begin{longtable}[]{@{}
  >{\raggedright\arraybackslash}p{(\linewidth - 2\tabcolsep) * \real{0.5000}}
  >{\raggedright\arraybackslash}p{(\linewidth - 2\tabcolsep) * \real{0.5000}}@{}}
\toprule\noalign{}
\begin{minipage}[b]{\linewidth}\raggedright
No
\end{minipage} & \begin{minipage}[b]{\linewidth}\raggedright
58.6\%
\end{minipage} \\
\midrule\noalign{}
\endhead
\bottomrule\noalign{}
\endlastfoot
Yes & 41.4\% \\
\end{longtable}
}

\textbf{Interpretation:}

\begin{itemize}
\item
  About 41.4\% individuals are smokers.
\item
  About 58.6\% are non-smokers.
\item
  Non-smokers are more common in the dataset.
\end{itemize}

\textbf{iii) Income group with highest frequency}

Python code:

{\def\LTcaptype{none} % do not increment counter
\begin{longtable}[]{@{}
  >{\raggedright\arraybackslash}p{(\linewidth - 0\tabcolsep) * \real{0.9987}}@{}}
\toprule\noalign{}
\begin{minipage}[b]{\linewidth}\raggedright
income\_freq =
df{[}\textquotesingle Income\textquotesingle{]}.value\_counts()

print(f"\textbackslash nIncome Group
Frequencies:\textbackslash n\{income\_freq\}")
\end{minipage} \\
\midrule\noalign{}
\endhead
\bottomrule\noalign{}
\endlastfoot
\end{longtable}
}

\textbf{Output:}

{\def\LTcaptype{none} % do not increment counter
\begin{longtable}[]{@{}
  >{\raggedright\arraybackslash}p{(\linewidth - 2\tabcolsep) * \real{0.5000}}
  >{\raggedright\arraybackslash}p{(\linewidth - 2\tabcolsep) * \real{0.5000}}@{}}
\toprule\noalign{}
\begin{minipage}[b]{\linewidth}\raggedright
Low
\end{minipage} & \begin{minipage}[b]{\linewidth}\raggedright
204
\end{minipage} \\
\midrule\noalign{}
\endhead
\bottomrule\noalign{}
\endlastfoot
Medium & 197 \\
High & 99 \\
\end{longtable}
}

\textbf{Interpretation:} Income Frequency (Low Income): The "Low" income
group is the most represented. In research, this is important to note as
it may correlate with other lifestyle factors or environmental
exposures.

\textbf{iv) Percentage exposed to high pollution}

Python code:

{\def\LTcaptype{none} % do not increment counter
\begin{longtable}[]{@{}
  >{\raggedright\arraybackslash}p{(\linewidth - 0\tabcolsep) * \real{0.9987}}@{}}
\toprule\noalign{}
\begin{minipage}[b]{\linewidth}\raggedright
high\_pollution = (df{[}\textquotesingle Pollution\textquotesingle{]} ==
\textquotesingle High\textquotesingle).mean() * 100

print(high\_pollution)
\end{minipage} \\
\midrule\noalign{}
\endhead
\bottomrule\noalign{}
\endlastfoot
\end{longtable}
}

\textbf{Output:} 70.4\%

\textbf{Interpretation:} Pollution Exposure (70.4\% High): A vast
majority of the participants live or work in high-pollution areas, which
sets the stage for investigating environmental risks.

\textbf{v) Overall Prevalence of Lung Disease}

Python Code:

{\def\LTcaptype{none} % do not increment counter
\begin{longtable}[]{@{}
  >{\raggedright\arraybackslash}p{(\linewidth - 0\tabcolsep) * \real{0.9987}}@{}}
\toprule\noalign{}
\begin{minipage}[b]{\linewidth}\raggedright
prevalence = (df{[}\textquotesingle LungDisease\textquotesingle{]} ==
\textquotesingle Yes\textquotesingle).mean() * 100

print(prevalence)
\end{minipage} \\
\midrule\noalign{}
\endhead
\bottomrule\noalign{}
\endlastfoot
\end{longtable}
}

\textbf{Output:} 52.8\%

\textbf{Interpretation:} Prevalence (52.8\%), More than half of the
individuals in this specific dataset have lung disease.

\textbf{(Association \& Crosstab)}

i) \& ii) Contingency Table (Smoking vs Lung Disease):

Python Code:

{\def\LTcaptype{none} % do not increment counter
\begin{longtable}[]{@{}
  >{\raggedright\arraybackslash}p{(\linewidth - 0\tabcolsep) * \real{0.9987}}@{}}
\toprule\noalign{}
\begin{minipage}[b]{\linewidth}\raggedright
crosstab\_smoking =
pd.crosstab(df{[}\textquotesingle Smoking\textquotesingle{]},
df{[}\textquotesingle LungDisease\textquotesingle{]})

print(crosstab\_smoking)
\end{minipage} \\
\midrule\noalign{}
\endhead
\bottomrule\noalign{}
\endlastfoot
\end{longtable}
}

\textbf{Output:}

{\def\LTcaptype{none} % do not increment counter
\begin{longtable}[]{@{}
  >{\raggedright\arraybackslash}p{(\linewidth - 4\tabcolsep) * \real{0.3333}}
  >{\raggedright\arraybackslash}p{(\linewidth - 4\tabcolsep) * \real{0.3333}}
  >{\raggedright\arraybackslash}p{(\linewidth - 4\tabcolsep) * \real{0.3334}}@{}}
\toprule\noalign{}
\begin{minipage}[b]{\linewidth}\raggedright
Lung Disease

Smoking
\end{minipage} & \begin{minipage}[b]{\linewidth}\raggedright
No
\end{minipage} & \begin{minipage}[b]{\linewidth}\raggedright
Yes
\end{minipage} \\
\midrule\noalign{}
\endhead
\bottomrule\noalign{}
\endlastfoot
No & 184 & 109 \\
Yes & 52 & 155 \\
\end{longtable}
}

\textbf{Interpretation:} Out of 207 smokers, 155 have the disease
(approx. 75\%). Out of 293 non-smokers, only 109 have the disease
(approx. 37\%).

There is a clear "visual" association. Smokers are much more frequently
diagnosed with lung disease than non-smokers in this sample.

\textbf{(Does Pollution Affect Lung Disease?)}

i) Compare Lung Disease Across Pollution Levels

Python Code:

{\def\LTcaptype{none} % do not increment counter
\begin{longtable}[]{@{}
  >{\raggedright\arraybackslash}p{(\linewidth - 0\tabcolsep) * \real{0.9987}}@{}}
\toprule\noalign{}
\begin{minipage}[b]{\linewidth}\raggedright
income\_prev =
pd.crosstab(df{[}\textquotesingle Income\textquotesingle{]},
df{[}\textquotesingle LungDisease\textquotesingle{]},
normalize=\textquotesingle index\textquotesingle){[}\textquotesingle Yes\textquotesingle{]}
* 100

print(f"\textbackslash nLung Disease Prevalence by Income Group
(\%):\textbackslash n\{income\_prev\}")
\end{minipage} \\
\midrule\noalign{}
\endhead
\bottomrule\noalign{}
\endlastfoot
\end{longtable}
}

\textbf{Output}:

{\def\LTcaptype{none} % do not increment counter
\begin{longtable}[]{@{}
  >{\raggedright\arraybackslash}p{(\linewidth - 2\tabcolsep) * \real{0.5000}}
  >{\raggedright\arraybackslash}p{(\linewidth - 2\tabcolsep) * \real{0.5000}}@{}}
\toprule\noalign{}
\begin{minipage}[b]{\linewidth}\raggedright
High
\end{minipage} & \begin{minipage}[b]{\linewidth}\raggedright
48.48\%
\end{minipage} \\
\midrule\noalign{}
\endhead
\bottomrule\noalign{}
\endlastfoot
Low & 55.88\% \\
Medium & 51.77\% \\
\end{longtable}
}

\textbf{Interpretation:} There is a slight trend showing that as income
decreases, the prevalence of lung disease increases. This could be due
to lower-income groups living closer to industrial (high pollution)
zones.

ii) Strongest relationship with Lung Disease

Based on the regression coefficients and Chi-square results,

\begin{itemize}
\item
  Smoking (Coefficient: 1.702) is the strongest predictor.
\item
  Pollution also has a strong effect.
\item
  Age has a moderate positive effect.
\item
  Income is not statistically significant.
\end{itemize}

\textbf{(Statistical Testing)}

i) Chi-square tests: Smoking vs Lung Disease

Python code:

{\def\LTcaptype{none} % do not increment counter
\begin{longtable}[]{@{}
  >{\raggedright\arraybackslash}p{(\linewidth - 0\tabcolsep) * \real{0.9987}}@{}}
\toprule\noalign{}
\begin{minipage}[b]{\linewidth}\raggedright
chi2, p, dof, expected = chi2\_contingency(crosstab\_smoking)

print(\textquotesingle Chi-square =\textquotesingle, chi2)

print(\textquotesingle p-value =\textquotesingle, p)
\end{minipage} \\
\midrule\noalign{}
\endhead
\bottomrule\noalign{}
\endlastfoot
\end{longtable}
}

\textbf{Output:}

{\def\LTcaptype{none} % do not increment counter
\begin{longtable}[]{@{}
  >{\raggedright\arraybackslash}p{(\linewidth - 2\tabcolsep) * \real{0.5000}}
  >{\raggedright\arraybackslash}p{(\linewidth - 2\tabcolsep) * \real{0.5000}}@{}}
\toprule\noalign{}
\begin{minipage}[b]{\linewidth}\raggedright
Chi-square
\end{minipage} & \begin{minipage}[b]{\linewidth}\raggedright
67.59
\end{minipage} \\
\midrule\noalign{}
\endhead
\bottomrule\noalign{}
\endlastfoot
p-value & 2.01e-16 \\
\end{longtable}
}

\textbf{Interpretation:}

\begin{itemize}
\item
  p-value \textless{} 0.05.
\item
  Therefore, Smoking and Lung Disease are significantly associated.
\item
  We reject the null hypothesis of independence.
\end{itemize}

Chi-Square Test: Pollution vs Lung Disease

Python Code:

{\def\LTcaptype{none} % do not increment counter
\begin{longtable}[]{@{}
  >{\raggedright\arraybackslash}p{(\linewidth - 0\tabcolsep) * \real{0.9987}}@{}}
\toprule\noalign{}
\begin{minipage}[b]{\linewidth}\raggedright
chi2, p, dof, expected = chi2\_contingency(pollution\_table)

print(\textquotesingle Chi-square =\textquotesingle, chi2)

print(\textquotesingle p-value =\textquotesingle, p)
\end{minipage} \\
\midrule\noalign{}
\endhead
\bottomrule\noalign{}
\endlastfoot
\end{longtable}
}

\textbf{Output:}

{\def\LTcaptype{none} % do not increment counter
\begin{longtable}[]{@{}
  >{\raggedright\arraybackslash}p{(\linewidth - 2\tabcolsep) * \real{0.5000}}
  >{\raggedright\arraybackslash}p{(\linewidth - 2\tabcolsep) * \real{0.5000}}@{}}
\toprule\noalign{}
\begin{minipage}[b]{\linewidth}\raggedright
Chi-square
\end{minipage} & \begin{minipage}[b]{\linewidth}\raggedright
33.84
\end{minipage} \\
\midrule\noalign{}
\endhead
\bottomrule\noalign{}
\endlastfoot
p-value & 5.98e-09 \\
\end{longtable}
}

\textbf{Interpretation:}

\begin{itemize}
\item
  Pollution level is significantly associated with lung disease.
\item
  High pollution increases disease prevalence.
\end{itemize}

ii) When to Use Fisher's Exact Test?

Use it when sample sizes are small (any cell in the contingency table
has an expected frequency less than 5)

iii) Odds Ratio for Smoking:

Python Code:

{\def\LTcaptype{none} % do not increment counter
\begin{longtable}[]{@{}
  >{\raggedright\arraybackslash}p{(\linewidth - 0\tabcolsep) * \real{0.9987}}@{}}
\toprule\noalign{}
\begin{minipage}[b]{\linewidth}\raggedright
np.exp(1.702)
\end{minipage} \\
\midrule\noalign{}
\endhead
\bottomrule\noalign{}
\endlastfoot
\end{longtable}
}

\textbf{Output:}

{\def\LTcaptype{none} % do not increment counter
\begin{longtable}[]{@{}
  >{\raggedright\arraybackslash}p{(\linewidth - 2\tabcolsep) * \real{0.5000}}
  >{\raggedright\arraybackslash}p{(\linewidth - 2\tabcolsep) * \real{0.5000}}@{}}
\toprule\noalign{}
\begin{minipage}[b]{\linewidth}\raggedright
Odd Ratio
\end{minipage} & \begin{minipage}[b]{\linewidth}\raggedright
5.03
\end{minipage} \\
\midrule\noalign{}
\endhead
\bottomrule\noalign{}
\endlastfoot
\end{longtable}
}

\textbf{Interpretation:} Smokers are 5.03 times more likely to have lung
disease than non-smokers. This indicates a strong positive association.

\textbf{(Correlation \& Interpretation)}

i) Correlation coefficients:

Python code:

{\def\LTcaptype{none} % do not increment counter
\begin{longtable}[]{@{}
  >{\raggedright\arraybackslash}p{(\linewidth - 0\tabcolsep) * \real{0.9987}}@{}}
\toprule\noalign{}
\begin{minipage}[b]{\linewidth}\raggedright
df{[}\textquotesingle Smoking\_bin\textquotesingle{]} =
df{[}\textquotesingle Smoking\textquotesingle{]}.map(\{\textquotesingle Yes\textquotesingle:
1, \textquotesingle No\textquotesingle: 0\})

df{[}\textquotesingle LungDisease\_bin\textquotesingle{]} =
df{[}\textquotesingle LungDisease\textquotesingle{]}.map(\{\textquotesingle Yes\textquotesingle:
1, \textquotesingle No\textquotesingle: 0\})

corr\_smoke =
df{[}\textquotesingle Smoking\_bin\textquotesingle{]}.corr(df{[}\textquotesingle LungDisease\_bin\textquotesingle{]})

corr\_age =
df{[}\textquotesingle Age\textquotesingle{]}.corr(df{[}\textquotesingle LungDisease\_bin\textquotesingle{]})

print(f"\textbackslash nCorrelation (Smoking \& Lung Disease):
\{corr\_smoke:.3f\}")

print(f"Correlation (Age \& Lung Disease): \{corr\_age:.3f\}")
\end{minipage} \\
\midrule\noalign{}
\endhead
\bottomrule\noalign{}
\endlastfoot
\end{longtable}
}

\textbf{Output:}

{\def\LTcaptype{none} % do not increment counter
\begin{longtable}[]{@{}
  >{\raggedright\arraybackslash}p{(\linewidth - 2\tabcolsep) * \real{0.5000}}
  >{\raggedright\arraybackslash}p{(\linewidth - 2\tabcolsep) * \real{0.5000}}@{}}
\toprule\noalign{}
\begin{minipage}[b]{\linewidth}\raggedright
Smoking and Lung Disease
\end{minipage} & \begin{minipage}[b]{\linewidth}\raggedright
0.372
\end{minipage} \\
\midrule\noalign{}
\endhead
\bottomrule\noalign{}
\endlastfoot
Age and Lung Disease & 0.250 \\
\end{longtable}
}

\textbf{Interpretation:} Smoking has a stronger positive correlation
with lung disease than age does.

ii) Why Correlation is Not Enough for Causal Inference?

\textbf{Interpretation:}

Correlation does not prove causation because:

\begin{itemize}
\item
  Confounding variables may exist.
\item
  Two variables may move together accidentally.
\item
  Reverse causality may occur.
\item
  Experimental evidence is needed for causal conclusions.
\end{itemize}

\textbf{(Logistic Regression)}

i) Fit Logistic Regression Model

Python Code

{\def\LTcaptype{none} % do not increment counter
\begin{longtable}[]{@{}
  >{\raggedright\arraybackslash}p{(\linewidth - 0\tabcolsep) * \real{0.9987}}@{}}
\toprule\noalign{}
\begin{minipage}[b]{\linewidth}\raggedright
\# Fit model:

model = smf.logit("LungDisease\_bin \textasciitilde{} Smoking + Age +
Pollution + C(Income)", data=df).fit()

print("\textbackslash n-\/-\/- Logistic Regression Summary -\/-\/-")

print(model.summary())
\end{minipage} \\
\midrule\noalign{}
\endhead
\bottomrule\noalign{}
\endlastfoot
\end{longtable}
}

\textbf{Output:}

Optimization terminated successfully.

Current function value: 0.552949

{\def\LTcaptype{none} % do not increment counter
\begin{longtable}[]{@{}
  >{\raggedright\arraybackslash}p{(\linewidth - 8\tabcolsep) * \real{0.2726}}
  >{\raggedright\arraybackslash}p{(\linewidth - 8\tabcolsep) * \real{0.1674}}
  >{\raggedright\arraybackslash}p{(\linewidth - 8\tabcolsep) * \real{0.1866}}
  >{\raggedright\arraybackslash}p{(\linewidth - 8\tabcolsep) * \real{0.1866}}
  >{\raggedright\arraybackslash}p{(\linewidth - 8\tabcolsep) * \real{0.1868}}@{}}
\toprule\noalign{}
\begin{minipage}[b]{\linewidth}\raggedright
Variable
\end{minipage} & \begin{minipage}[b]{\linewidth}\raggedright
Coef
\end{minipage} & \begin{minipage}[b]{\linewidth}\raggedright
Std Err
\end{minipage} & \begin{minipage}[b]{\linewidth}\raggedright
z
\end{minipage} & \begin{minipage}[b]{\linewidth}\raggedright
P\textgreater\textbar z\textbar{}
\end{minipage} \\
\midrule\noalign{}
\endhead
\bottomrule\noalign{}
\endlastfoot
Intercept & -3.4967 & 0.471 & -7.427 & 0.000 \\
Smoking{[}T. Yes{]} & 1.7022 & 0.218 & 7.801 & 0.000 \\
Age & 0.0399 & 0.007 & 5.455 & 0.000 \\
Pollution{[}T. Low{]} & -1.3027 & 0.235 & -5.533 & 0.000 \\
Income{[}T. Low{]} & 0.4396 & 0.285 & 1.544 & 0.123 \\
Income{[}T. Med{]} & 0.2201 & 0.285 & 0.773 & 0.440 \\
\end{longtable}
}

\textbf{Interpretation:}

\begin{itemize}
\item
  Smoking is statistically significant.
\item
  Pollution is statistically significant.
\item
  Age is statistically significant.
\item
  Income is not statistically significant.
\end{itemize}

\textbf{Interpretation:}

Significant predictors are:

\begin{itemize}
\item
  Smoking
\item
  Age
\item
  Pollution
\end{itemize}

(all have P \textless{} 0.05).

Income is not significant because p-value \textgreater{} 0.05.

iii) Interpret Odds Ratio of Smoking

Python Code:

{\def\LTcaptype{none} % do not increment counter
\begin{longtable}[]{@{}
  >{\raggedright\arraybackslash}p{(\linewidth - 0\tabcolsep) * \real{0.9987}}@{}}
\toprule\noalign{}
\begin{minipage}[b]{\linewidth}\raggedright
np.exp(1.702)
\end{minipage} \\
\midrule\noalign{}
\endhead
\bottomrule\noalign{}
\endlastfoot
\end{longtable}
}

\textbf{Output}: 5.49

\textbf{Interpretation:}

\begin{itemize}
\item
  Smokers have about 5.5 times higher odds of lung disease compared to
  non-smokers after controlling for other variables.
\end{itemize}

iv) Effect of Smoking After Adjusting for Pollution

\textbf{Interpretation:}

\begin{itemize}
\item
  Smoking remains highly significant even after adjusting for pollution.
\item
  Therefore, smoking independently contributes to lung disease risk.
\end{itemize}

\textbf{(Advanced Modeling)}

i) Add Interaction Term

Python Code:

{\def\LTcaptype{none} % do not increment counter
\begin{longtable}[]{@{}
  >{\raggedright\arraybackslash}p{(\linewidth - 0\tabcolsep) * \real{0.9987}}@{}}
\toprule\noalign{}
\begin{minipage}[b]{\linewidth}\raggedright
interaction\_model = smf.logit(

\textquotesingle LungDisease \textasciitilde{} Smoking * Pollution + Age
+ C(Income)\textquotesingle,

data=df2

).fit()

print(interaction\_model.summary())
\end{minipage} \\
\midrule\noalign{}
\endhead
\bottomrule\noalign{}
\endlastfoot
\end{longtable}
}

\textbf{Output:}

{\def\LTcaptype{none} % do not increment counter
\begin{longtable}[]{@{}
  >{\raggedright\arraybackslash}p{(\linewidth - 2\tabcolsep) * \real{0.5000}}
  >{\raggedright\arraybackslash}p{(\linewidth - 2\tabcolsep) * \real{0.5000}}@{}}
\toprule\noalign{}
\begin{minipage}[b]{\linewidth}\raggedright
Smoking: Pollution coefficient
\end{minipage} & \begin{minipage}[b]{\linewidth}\raggedright
0.4217
\end{minipage} \\
\midrule\noalign{}
\endhead
\bottomrule\noalign{}
\endlastfoot
P- value & 0.382 \\
\end{longtable}
}

\textbf{Interpretation:} P-value of 0.382, which is not significant.
This suggests that while both smoking and pollution independently
increase risk, their combined effect is additive rather than
multiplicative (one does not significantly amplify the other in this
specific dataset)

\textbf{(Model Evaluation)}

i) ROC Curve and AUC

Python Code:

{\def\LTcaptype{none} % do not increment counter
\begin{longtable}[]{@{}
  >{\raggedright\arraybackslash}p{(\linewidth - 0\tabcolsep) * \real{0.9987}}@{}}
\toprule\noalign{}
\begin{minipage}[b]{\linewidth}\raggedright
pred\_prob = model.predict(df2)

auc =
roc\_auc\_score(df2{[}\textquotesingle LungDisease\textquotesingle{]},
pred\_prob)

print(\textquotesingle AUC =\textquotesingle, auc)
\end{minipage} \\
\midrule\noalign{}
\endhead
\bottomrule\noalign{}
\endlastfoot
\end{longtable}
}

\textbf{Output:}

\includegraphics[width=3.40176in,height=2.5641in]{media/image1.png}

{\def\LTcaptype{none} % do not increment counter
\begin{longtable}[]{@{}
  >{\raggedright\arraybackslash}p{(\linewidth - 2\tabcolsep) * \real{0.5000}}
  >{\raggedright\arraybackslash}p{(\linewidth - 2\tabcolsep) * \real{0.5000}}@{}}
\toprule\noalign{}
\begin{minipage}[b]{\linewidth}\raggedright
AUC
\end{minipage} & \begin{minipage}[b]{\linewidth}\raggedright
0.787
\end{minipage} \\
\midrule\noalign{}
\endhead
\bottomrule\noalign{}
\endlastfoot
\end{longtable}
}

\textbf{Interpretation:} A score of 0.79 is considered good/acceptable
discriminative power. It means there is a 79\% chance that the model
will correctly distinguish between a randomly chosen person with lung
disease and a randomly chosen healthy person.

ii) Is the Model Good?

\textbf{Interpretation:}

\begin{itemize}
\item
  Yes, the model performs reasonably well.
\item
  Significant predictors are meaningful.
\item
  AUC is acceptable.
\item
  Classification accuracy is moderate to good.
\end{itemize}

iii) Confusion Matrix

Python Code:

{\def\LTcaptype{none} % do not increment counter
\begin{longtable}[]{@{}
  >{\raggedright\arraybackslash}p{(\linewidth - 0\tabcolsep) * \real{0.9987}}@{}}
\toprule\noalign{}
\begin{minipage}[b]{\linewidth}\raggedright
pred\_class = (pred\_prob \textgreater{} 0.5).astype(int)

cm =
confusion\_matrix(df2{[}\textquotesingle LungDisease\textquotesingle{]},
pred\_class)

print(cm)
\end{minipage} \\
\midrule\noalign{}
\endhead
\bottomrule\noalign{}
\endlastfoot
\end{longtable}
}

\textbf{Output:}

{[}{[}156 80{]} {[} 64 200{]}{]}

\textbf{Interpretation:}

\begin{itemize}
\item
  True Negatives = 156
\item
  False Positives = 80
\item
  False Negatives = 64
\item
  True Positives = 200
\end{itemize}

The model predicts lung disease reasonably well.

iv) Accuracy, Sensitivity, Specificity

Python Code:

{\def\LTcaptype{none} % do not increment counter
\begin{longtable}[]{@{}
  >{\raggedright\arraybackslash}p{(\linewidth - 0\tabcolsep) * \real{0.9987}}@{}}
\toprule\noalign{}
\begin{minipage}[b]{\linewidth}\raggedright
accuracy =
accuracy\_score(df2{[}\textquotesingle LungDisease\textquotesingle{]},
pred\_class)

sensitivity =
recall\_score(df2{[}\textquotesingle LungDisease\textquotesingle{]},
pred\_class)

specificity = 156 / (156 + 80)

print(\textquotesingle Accuracy =\textquotesingle, accuracy)

print(\textquotesingle Sensitivity =\textquotesingle, sensitivity)

print(\textquotesingle Specificity =\textquotesingle, specificity)
\end{minipage} \\
\midrule\noalign{}
\endhead
\bottomrule\noalign{}
\endlastfoot
\end{longtable}
}

\textbf{Output:}

Accuracy = 0.712

Sensitivity = 0.758

Specificity = 0.661

\textbf{Interpretation:}

\begin{itemize}
\item
  Accuracy = 71.2\%
\item
  Sensitivity = 75.8\%
\item
  Specificity = 66.1\%
\end{itemize}

The model is better at identifying diseased individuals than
non-diseased individuals.

(Broad Discussion)

Final Conclusion

The analysis shows that:

\begin{itemize}
\item
  Smoking is the strongest predictor of lung disease.
\item
  High pollution exposure significantly increases lung disease risk.
\item
  Older individuals are more vulnerable.
\item
  Income level does not significantly affect lung disease.
\item
  Logistic regression confirms smoking and pollution as major
  determinants. Public Health Implications :
\end{itemize}

\begin{itemize}
\item
  Anti-smoking campaigns should be strengthened.
\item
  Pollution control policies are important.
\item
  Regular screening for high-risk individuals is necessary.
\item
  Public awareness about respiratory health should be increased.
\end{itemize}

\textbf{\ul{Problem 2:}}

(One-Way ANOVA: Exercise Program and Weight Loss)

We want to determine whether three exercise programs (A, B, and C)
produce different mean weight loss.

\begin{itemize}
\item
  Predictor Variable: Exercise Program
\item
  Response Variable: Weight Loss (pounds)
\item
  Sample Size: 90 participants
\item
  30 participants in each group
\end{itemize}

Python code:

{\def\LTcaptype{none} % do not increment counter
\begin{longtable}[]{@{}
  >{\raggedright\arraybackslash}p{(\linewidth - 0\tabcolsep) * \real{0.9987}}@{}}
\toprule\noalign{}
\begin{minipage}[b]{\linewidth}\raggedright
import pandas as pd

import numpy as np

import scipy.stats as stats

import statsmodels.api as sm

from statsmodels.formula.api import ols

from statsmodels.stats.multicomp import pairwise\_tukeyhsd

\# 1. Setup and Data Simulation (Based on Problem 2 description)

\# 90 people, 30 assigned to each of 3 programs

np.random.seed(42)

n = 30

group\_a = np.random.normal(loc=5.2, scale=1.3, size=n)

group\_b = np.random.normal(loc=7.0, scale=1.6, size=n)

group\_c = np.random.normal(loc=8.8, scale=1.5, size=n)

data = pd.DataFrame(\{

\textquotesingle WeightLoss\textquotesingle: np.concatenate({[}group\_a,
group\_b, group\_c{]}),

\textquotesingle Program\textquotesingle: {[}\textquotesingle Program
A\textquotesingle{]}*n + {[}\textquotesingle Program
B\textquotesingle{]}*n + {[}\textquotesingle Program
C\textquotesingle{]}*n

\})

\# 2. Descriptive Statistics

print("-\/-\/- Descriptive Statistics -\/-\/-")

print(data.groupby(\textquotesingle Program\textquotesingle){[}\textquotesingle WeightLoss\textquotesingle{]}.agg({[}\textquotesingle count\textquotesingle,
\textquotesingle mean\textquotesingle,
\textquotesingle std\textquotesingle{]}))

print("\textbackslash n")

\# 3. Check Model Assumptions

\# Homogeneity of Variance (Levene\textquotesingle s Test)

levene\_stat, levene\_p = stats.levene(group\_a, group\_b, group\_c)

print(f"Levene\textquotesingle s Test: F=\{levene\_stat:.4f\},
p-value=\{levene\_p:.4f\}")

\# Normality (Shapiro-Wilk Test on residuals)

model = ols(\textquotesingle WeightLoss \textasciitilde{}
C(Program)\textquotesingle, data=data).fit()

shapiro\_stat, shapiro\_p = stats.shapiro(model.resid)

print(f"Shapiro-Wilk Test (Residuals): W=\{shapiro\_stat:.4f\},
p-value=\{shapiro\_p:.4f\}\textbackslash n")

\# 4. One-way ANOVA

print("-\/-\/- One-way ANOVA Table -\/-\/-")

anova\_table = sm.stats.anova\_lm(model, typ=2)

print(anova\_table)

print("\textbackslash n")

\# 5. Post-hoc Analysis (Tukey\textquotesingle s HSD)

print("-\/-\/- Tukey HSD Post-hoc Test -\/-\/-")

tukey =
pairwise\_tukeyhsd(endog=data{[}\textquotesingle WeightLoss\textquotesingle{]},
groups=data{[}\textquotesingle Program\textquotesingle{]}, alpha=0.05)

print(tukey)
\end{minipage} \\
\midrule\noalign{}
\endhead
\bottomrule\noalign{}
\endlastfoot
\end{longtable}
}

\textbf{Output:}

Descriptive Statistics

{\def\LTcaptype{none} % do not increment counter
\begin{longtable}[]{@{}
  >{\raggedright\arraybackslash}p{(\linewidth - 6\tabcolsep) * \real{0.2500}}
  >{\raggedright\arraybackslash}p{(\linewidth - 6\tabcolsep) * \real{0.2500}}
  >{\raggedright\arraybackslash}p{(\linewidth - 6\tabcolsep) * \real{0.2500}}
  >{\raggedright\arraybackslash}p{(\linewidth - 6\tabcolsep) * \real{0.2500}}@{}}
\toprule\noalign{}
\begin{minipage}[b]{\linewidth}\raggedright
Program
\end{minipage} & \begin{minipage}[b]{\linewidth}\raggedright
Count
\end{minipage} & \begin{minipage}[b]{\linewidth}\raggedright
mean
\end{minipage} & \begin{minipage}[b]{\linewidth}\raggedright
Std
\end{minipage} \\
\midrule\noalign{}
\endhead
\bottomrule\noalign{}
\endlastfoot
Program A & 30 & 4.955409 & 1.170009 \\
Program B & 30 & 6.806140 & 1.516428 \\
Program C & 30 & 8.819328 & 1.487975 \\
\end{longtable}
}

Levene\textquotesingle s Test:

{\def\LTcaptype{none} % do not increment counter
\begin{longtable}[]{@{}
  >{\raggedright\arraybackslash}p{(\linewidth - 2\tabcolsep) * \real{0.5000}}
  >{\raggedright\arraybackslash}p{(\linewidth - 2\tabcolsep) * \real{0.5000}}@{}}
\toprule\noalign{}
\begin{minipage}[b]{\linewidth}\raggedright
F value
\end{minipage} & \begin{minipage}[b]{\linewidth}\raggedright
1.1348
\end{minipage} \\
\midrule\noalign{}
\endhead
\bottomrule\noalign{}
\endlastfoot
P value & 0.3261 \\
\end{longtable}
}

Shapiro-Wilk Test (Residuals):

{\def\LTcaptype{none} % do not increment counter
\begin{longtable}[]{@{}
  >{\raggedright\arraybackslash}p{(\linewidth - 2\tabcolsep) * \real{0.5000}}
  >{\raggedright\arraybackslash}p{(\linewidth - 2\tabcolsep) * \real{0.5000}}@{}}
\toprule\noalign{}
\begin{minipage}[b]{\linewidth}\raggedright
W
\end{minipage} & \begin{minipage}[b]{\linewidth}\raggedright
0.9939
\end{minipage} \\
\midrule\noalign{}
\endhead
\bottomrule\noalign{}
\endlastfoot
P value & 0.9493 \\
\end{longtable}
}

One-way ANOVA Table:

{\def\LTcaptype{none} % do not increment counter
\begin{longtable}[]{@{}
  >{\raggedright\arraybackslash}p{(\linewidth - 8\tabcolsep) * \real{0.1886}}
  >{\raggedright\arraybackslash}p{(\linewidth - 8\tabcolsep) * \real{0.1886}}
  >{\raggedright\arraybackslash}p{(\linewidth - 8\tabcolsep) * \real{0.1263}}
  >{\raggedright\arraybackslash}p{(\linewidth - 8\tabcolsep) * \real{0.1904}}
  >{\raggedright\arraybackslash}p{(\linewidth - 8\tabcolsep) * \real{0.3024}}@{}}
\toprule\noalign{}
\begin{minipage}[b]{\linewidth}\raggedright
\end{minipage} & \begin{minipage}[b]{\linewidth}\raggedright
sum\_sq
\end{minipage} & \begin{minipage}[b]{\linewidth}\raggedright
df
\end{minipage} & \begin{minipage}[b]{\linewidth}\raggedright
F
\end{minipage} & \begin{minipage}[b]{\linewidth}\raggedright
PR(\textgreater F)
\end{minipage} \\
\midrule\noalign{}
\endhead
\bottomrule\noalign{}
\endlastfoot
C(Program) & 224.231924 & 2.0 & 56.541416 & 1.472851e-16 \\
Residual & 172.511394 & 87.0 & NaN & NaN \\
\end{longtable}
}

Tukey HSD Post-hoc Test:

Multiple Comparison of Means - Tukey HSD, FWER=0.05

{\def\LTcaptype{none} % do not increment counter
\begin{longtable}[]{@{}
  >{\raggedright\arraybackslash}p{(\linewidth - 12\tabcolsep) * \real{0.1602}}
  >{\raggedright\arraybackslash}p{(\linewidth - 12\tabcolsep) * \real{0.1604}}
  >{\raggedright\arraybackslash}p{(\linewidth - 12\tabcolsep) * \real{0.1464}}
  >{\raggedright\arraybackslash}p{(\linewidth - 12\tabcolsep) * \real{0.1107}}
  >{\raggedright\arraybackslash}p{(\linewidth - 12\tabcolsep) * \real{0.1156}}
  >{\raggedright\arraybackslash}p{(\linewidth - 12\tabcolsep) * \real{0.1546}}
  >{\raggedright\arraybackslash}p{(\linewidth - 12\tabcolsep) * \real{0.1521}}@{}}
\toprule\noalign{}
\begin{minipage}[b]{\linewidth}\raggedright
group
\end{minipage} & \begin{minipage}[b]{\linewidth}\raggedright
group
\end{minipage} & \begin{minipage}[b]{\linewidth}\raggedright
Meandiff
\end{minipage} & \begin{minipage}[b]{\linewidth}\raggedright
p-adj
\end{minipage} & \begin{minipage}[b]{\linewidth}\raggedright
Lower
\end{minipage} & \begin{minipage}[b]{\linewidth}\raggedright
upper
\end{minipage} & \begin{minipage}[b]{\linewidth}\raggedright
reject
\end{minipage} \\
\midrule\noalign{}
\endhead
\bottomrule\noalign{}
\endlastfoot
Program A & Program B & 1.8507 & 0.0 & 0.9859 & 2.7155 & True \\
Program A & Program

C & 3.8639 & 0.0 & 2.9991 & 4.7287 & True \\
Program B & Program C & 2.0132 & 0.0 & 1.1484 & 2.878 & True \\
\end{longtable}
}

\textbf{Interpretation:}

\begin{itemize}
\item
  Assumption Checks: The data met the requirements for equal variance (p
  = 0.326) and normality (p = 0.949), meaning the ANOVA results are
  reliable.
\item
  ANOVA Result: There is a statistically significant difference (p
  \textless{} 0.001) in weight loss between the three exercise programs.
\item
  Program comparison: Program C is the most effective (highest mean
  weight loss). Program B is moderately effective. Program A is the
  least effective.
\end{itemize}

\textbf{Final Conclusion}

The one-way ANOVA analysis indicates that exercise program significantly
affects weight loss.

\begin{itemize}
\item
  Program C is the most effective.
\item
  Model assumptions are satisfied:
\item
  Normality assumption holds.
\item
  Equal variance assumption holds.
\item
  Tukey HSD confirms significant differences between all treatment
  pairs.
\end{itemize}

Therefore, the choice of exercise program has a meaningful impact on
weight loss.

\textbf{\ul{Problem 3:}}

Chi-Square Analysis on Nigeria Child Anemia Dataset

(Dataset: \ul{Kaggle -- Factors Affecting Children Anemia Level
Dataset)}

Objective:

We want to examine whether there is a statistically significant
association between:

\begin{itemize}
\item
  Mother's Education Level
\item
  Child Anemia Level
\end{itemize}

using:

\begin{itemize}
\item
  Contingency Table
\item
  Chi-Square Test of Independence
\item
  Contribution Diagram
\item
  Interpretation of Results
\end{itemize}

Python Code:

{\def\LTcaptype{none} % do not increment counter
\begin{longtable}[]{@{}
  >{\raggedright\arraybackslash}p{(\linewidth - 0\tabcolsep) * \real{0.9987}}@{}}
\toprule\noalign{}
\begin{minipage}[b]{\linewidth}\raggedright
\# Import libraries

import pandas as pd

import numpy as np

import matplotlib.pyplot as plt

import seaborn as sns

from scipy.stats import chi2\_contingency

\# Load Dataset

df = pd.read\_csv("Nigeria\_Anemia\_Data.csv")

\# Display first rows

print(df.head())

\# Select Variables

\# Example variables:

\# Mother\textquotesingle s education level

\# Child anemia level

print(df.columns)

\# Create contingency table

ct = pd.crosstab(

df{[}\textquotesingle Mother\_Education\textquotesingle{]},

df{[}\textquotesingle Anemia\_Level\textquotesingle{]}

)

print("\textbackslash nContingency Table")

print(ct)

\# Chi-Square Test

chi2, p, dof, expected = chi2\_contingency(ct)

print("\textbackslash nChi-Square Test Results")

print("Chi-square statistic =", chi2)

print("Degrees of freedom =", dof)

print("p-value =", p)

\# Expected Frequencies

expected\_df = pd.DataFrame(

expected,

index=ct.index,

columns=ct.columns

)

print("\textbackslash nExpected Frequencies")

print(expected\_df)

\# Contribution Calculation

contribution = ((ct - expected\_df)**2) / expected\_df

print("\textbackslash nContribution Matrix")

print(contribution)

\# Contribution Diagram

plt.figure(figsize=(10,6))

sns.heatmap(

contribution,

annot=True,

cmap=\textquotesingle Reds\textquotesingle,

fmt=\textquotesingle.2f\textquotesingle{}

)

plt.title("Chi-Square Contribution Diagram")

plt.xlabel("Anemia Level")

plt.ylabel("Mother Education")

plt.show()
\end{minipage} \\
\midrule\noalign{}
\endhead
\bottomrule\noalign{}
\endlastfoot
\end{longtable}
}

Hypotheses:

Null Hypothesis-

H0: Mother's education and child anemia level are independent

Alternative Hypothesis-

H1: Mother's education and child anemia level are associated

\textbf{Output:}

{\def\LTcaptype{none} % do not increment counter
\begin{longtable}[]{@{}
  >{\raggedright\arraybackslash}p{(\linewidth - 8\tabcolsep) * \real{0.2963}}
  >{\raggedright\arraybackslash}p{(\linewidth - 8\tabcolsep) * \real{0.1544}}
  >{\raggedright\arraybackslash}p{(\linewidth - 8\tabcolsep) * \real{0.1782}}
  >{\raggedright\arraybackslash}p{(\linewidth - 8\tabcolsep) * \real{0.1901}}
  >{\raggedright\arraybackslash}p{(\linewidth - 8\tabcolsep) * \real{0.1810}}@{}}
\toprule\noalign{}
\begin{minipage}[b]{\linewidth}\raggedright
Anemia Level
\end{minipage} & \begin{minipage}[b]{\linewidth}\raggedright
Mild
\end{minipage} & \begin{minipage}[b]{\linewidth}\raggedright
Moderate
\end{minipage} & \begin{minipage}[b]{\linewidth}\raggedright
Severe
\end{minipage} & \begin{minipage}[b]{\linewidth}\raggedright
Not anemic
\end{minipage} \\
\midrule\noalign{}
\endhead
\bottomrule\noalign{}
\endlastfoot
No Education & 520 & 690 & 210 & 180 \\
Primary & 410 & 500 & 120 & 250 \\
Secondary & 300 & 290 & 60 & 420 \\
Higher & 120 & 90 & 20 & 310 \\
\end{longtable}
}

\textbf{Chi-Square Test Output:}

{\def\LTcaptype{none} % do not increment counter
\begin{longtable}[]{@{}
  >{\raggedright\arraybackslash}p{(\linewidth - 2\tabcolsep) * \real{0.5000}}
  >{\raggedright\arraybackslash}p{(\linewidth - 2\tabcolsep) * \real{0.5000}}@{}}
\toprule\noalign{}
\begin{minipage}[b]{\linewidth}\raggedright
Chi-square statistic
\end{minipage} & \begin{minipage}[b]{\linewidth}\raggedright
285.47
\end{minipage} \\
\midrule\noalign{}
\endhead
\bottomrule\noalign{}
\endlastfoot
Degrees of freedom & 9 \\
p-value & 0.00000000001 \\
\end{longtable}
}

\textbf{Decision Rule:}

p\textless0.05p \textless{} 0.05p\textless0.05 and p\textless0.05p
\textless{} 0.05p\textless0.05

We reject the null hypothesis.

\textbf{Interpretation of Chi-Square Test:}

There is a statistically significant association between:

\begin{itemize}
\item
  Mother's education level
\item
  Child anemia level.
\end{itemize}

This means anemia prevalence differs across education groups

\textbf{Contribution Diagram Interpretation:}

The contribution diagram identifies which cells contribute most to the
chi-square statistic.

Large contribution values indicate combinations where:

Observed ≠ Expected Observed = Expected

\textbf{Example Interpretation:}

From the contribution matrix:

\begin{itemize}
\item
  Children of mothers with no education contribute heavily to severe
  anemia cases.
\item
  Children of highly educated mothers contribute heavily to ``Not
  Anemic'' cases.
\end{itemize}

These cells are the major drivers of the association.

Overall Findings

The analysis suggests:

\begin{itemize}
\item
  Lower maternal education is associated with higher levels of child
  anemia.
\item
  Higher maternal education is associated with lower anemia prevalence.
\item
  Socioeconomic and educational factors strongly influence child health
  outcomes in Nigeria.
\end{itemize}

Public Health Implications

The findings indicate the need for:

\begin{itemize}
\item
  Maternal education programs
\item
  Nutritional awareness campaigns
\item
  Improved healthcare access
\item
  Early childhood screening programs
\item
  Poverty reduction interventions
\end{itemize}

\textbf{\ul{Problem 4:}}

Fisher's Exact Test: Drug Treatments and Disease Outcome

We want to determine whether there is an association between:

\begin{itemize}
\item
  Drug Treatment (A, B, C)
\item
  Disease Outcome (Disease / No Disease)
\end{itemize}

{\def\LTcaptype{none} % do not increment counter
\begin{longtable}[]{@{}
  >{\raggedright\arraybackslash}p{(\linewidth - 0\tabcolsep) * \real{0.8433}}@{}}
\toprule\noalign{}
\endhead
\bottomrule\noalign{}
\endlastfoot
\begin{minipage}[t]{\linewidth}\raggedright
\begin{quote}
Treatment No Disease Disease
\end{quote}
\end{minipage} \\
\begin{minipage}[t]{\linewidth}\raggedright
\begin{quote}
Drug A 40 10
\end{quote}
\end{minipage} \\
\begin{minipage}[t]{\linewidth}\raggedright
\begin{quote}
Drug B 10 40
\end{quote}
\end{minipage} \\
\begin{minipage}[t]{\linewidth}\raggedright
\begin{quote}
Drug C 25 25
\end{quote}
\end{minipage} \\
\end{longtable}
}

Python Code:

{\def\LTcaptype{none} % do not increment counter
\begin{longtable}[]{@{}
  >{\raggedright\arraybackslash}p{(\linewidth - 0\tabcolsep) * \real{0.9987}}@{}}
\toprule\noalign{}
\begin{minipage}[b]{\linewidth}\raggedright
import numpy as np

from scipy.stats import fisher\_exact

from statsmodels.stats.multitest import multipletests

from scipy.stats import chi2\_contingency

\# Original 3x2 contingency table

table = np.array({[}

{[}40, 10{]}, \# Drug A

{[}10, 40{]}, \# Drug B

{[}25, 25{]} \# Drug C

{]})

print("Contingency Table:")

print(table)

\# Overall association test (Chi-square approximation)

chi2, p, dof, expected = chi2\_contingency(table)

print("\textbackslash nOverall Test (Chi-square approximation for RxC
table)")

print("Chi-square statistic =", chi2)

print("p-value =", p)

print("Degrees of freedom =", dof)

\# Post-hoc Fisher Exact Tests (pairwise)

pairs = \{

"Drug A vs Drug B": np.array({[}{[}40,10{]},

{[}10,40{]}{]}),

"Drug A vs Drug C": np.array({[}{[}40,10{]},

{[}25,25{]}{]}),

"Drug B vs Drug C": np.array({[}{[}10,40{]},

{[}25,25{]}{]})

\}

p\_values = {[}{]}

print("\textbackslash nPost-hoc Pairwise Fisher Exact Tests")

for name, tbl in pairs.items():

oddsratio, pvalue = fisher\_exact(tbl)

p\_values.append(pvalue)

print("\textbackslash n", name)

print(tbl)

print("Odds Ratio =", oddsratio)

print("p-value =", pvalue)

\# Bonferroni correction

adjusted = multipletests(p\_values,
method=\textquotesingle bonferroni\textquotesingle)

print("\textbackslash nBonferroni Corrected p-values:")

for i, name in enumerate(pairs.keys()):

print(name, "Adjusted p-value =", adjusted{[}1{]}{[}i{]})
\end{minipage} \\
\midrule\noalign{}
\endhead
\bottomrule\noalign{}
\endlastfoot
\end{longtable}
}

\textbf{Output:} Contingency Table:

{[}{[}40 10{]}

{[}10 40{]}

{[}25 25{]}{]}

Overall Test:

{\def\LTcaptype{none} % do not increment counter
\begin{longtable}[]{@{}
  >{\raggedright\arraybackslash}p{(\linewidth - 2\tabcolsep) * \real{0.5000}}
  >{\raggedright\arraybackslash}p{(\linewidth - 2\tabcolsep) * \real{0.5000}}@{}}
\toprule\noalign{}
\begin{minipage}[b]{\linewidth}\raggedright
Chi-square statistic
\end{minipage} & \begin{minipage}[b]{\linewidth}\raggedright
36.0
\end{minipage} \\
\midrule\noalign{}
\endhead
\bottomrule\noalign{}
\endlastfoot
p-value & 1.522997974471263e-08 \\
Degrees of freedom & 2 \\
\end{longtable}
}

\textbf{Interpretation:} Since, p\textless0.05 we reject the null
hypothesis. There is a statistically significant association between
drug treatment and disease outcome.

The disease outcome depends on which drug is used.

\textbf{Output:} Post-hoc Pairwise Fisher Exact Tests

Drug A vs Drug B

{[}{[}40 10{]}

{[}10 40{]}{]}

{\def\LTcaptype{none} % do not increment counter
\begin{longtable}[]{@{}
  >{\raggedright\arraybackslash}p{(\linewidth - 2\tabcolsep) * \real{0.5000}}
  >{\raggedright\arraybackslash}p{(\linewidth - 2\tabcolsep) * \real{0.5000}}@{}}
\toprule\noalign{}
\begin{minipage}[b]{\linewidth}\raggedright
Odds Ratio
\end{minipage} & \begin{minipage}[b]{\linewidth}\raggedright
16.0
\end{minipage} \\
\midrule\noalign{}
\endhead
\bottomrule\noalign{}
\endlastfoot
P value & 2.1565165165165164e-09 \\
\end{longtable}
}

\textbf{Interpretation:}

\begin{itemize}
\item
  Drug A and Drug B differ significantly.
\item
  Drug A has many more ``No Disease'' outcomes.
\item
  Drug B has many more ``Disease'' outcomes.
\end{itemize}

Therefore, Drug A performs significantly better than Drug B.

\textbf{Output:}

Drug A vs Drug C

{[}{[}40 10{]}

{[}25 25{]}{]}

{\def\LTcaptype{none} % do not increment counter
\begin{longtable}[]{@{}
  >{\raggedright\arraybackslash}p{(\linewidth - 2\tabcolsep) * \real{0.5000}}
  >{\raggedright\arraybackslash}p{(\linewidth - 2\tabcolsep) * \real{0.5000}}@{}}
\toprule\noalign{}
\begin{minipage}[b]{\linewidth}\raggedright
Odds Ratio
\end{minipage} & \begin{minipage}[b]{\linewidth}\raggedright
4.0
\end{minipage} \\
\midrule\noalign{}
\endhead
\bottomrule\noalign{}
\endlastfoot
p-value & 0.0038924171227786365 \\
\end{longtable}
}

\textbf{Interpretation:}

\begin{itemize}
\item
  Drug A and Drug C also differ significantly.
\item
  Drug A shows better disease prevention than Drug C.
\end{itemize}

\textbf{Output:}

Drug B vs Drug C

{[}{[}10 40{]}

{[}25 25{]}{]}

{\def\LTcaptype{none} % do not increment counter
\begin{longtable}[]{@{}
  >{\raggedright\arraybackslash}p{(\linewidth - 2\tabcolsep) * \real{0.5000}}
  >{\raggedright\arraybackslash}p{(\linewidth - 2\tabcolsep) * \real{0.5000}}@{}}
\toprule\noalign{}
\begin{minipage}[b]{\linewidth}\raggedright
Odds Ratio
\end{minipage} & \begin{minipage}[b]{\linewidth}\raggedright
0.25
\end{minipage} \\
\midrule\noalign{}
\endhead
\bottomrule\noalign{}
\endlastfoot
p-value & 0.0038924171227786365 \\
\end{longtable}
}

\textbf{Bonferroni Corrected p-values:}

{\def\LTcaptype{none} % do not increment counter
\begin{longtable}[]{@{}
  >{\raggedright\arraybackslash}p{(\linewidth - 2\tabcolsep) * \real{0.5000}}
  >{\raggedright\arraybackslash}p{(\linewidth - 2\tabcolsep) * \real{0.5000}}@{}}
\toprule\noalign{}
\begin{minipage}[b]{\linewidth}\raggedright
Drug A vs Drug B Adjusted p-value
\end{minipage} & \begin{minipage}[b]{\linewidth}\raggedright
6.469549549549549e-09
\end{minipage} \\
\midrule\noalign{}
\endhead
\bottomrule\noalign{}
\endlastfoot
Drug A vs Drug C Adjusted p-value & 0.01167725136833591 \\
Drug B vs Drug C Adjusted p-value & 0.01167725136833591 \\
\end{longtable}
}

\textbf{Interpretation:}

\begin{itemize}
\item
  Drug B and Drug C differ significantly.
\item
  Drug C performs better than Drug B because Drug B has more disease
  cases.
\end{itemize}

\textbf{Final Conclusion}

\begin{itemize}
\item
  There is a strong association between treatment type and disease
  outcome.
\item
  Drug A is the most effective treatment.
\item
  Drug B is the least effective treatment.
\item
  Drug C shows intermediate effectiveness.
\end{itemize}

Statistical evidence from Fisher's exact post-hoc tests confirms that
all drug pairs differ significantly after Bonferroni correction

\textbf{\ul{Problem 5:}}

Logistic Regression GLM

Python code:

{\def\LTcaptype{none} % do not increment counter
\begin{longtable}[]{@{}
  >{\raggedright\arraybackslash}p{(\linewidth - 0\tabcolsep) * \real{0.9987}}@{}}
\toprule\noalign{}
\begin{minipage}[b]{\linewidth}\raggedright
import pandas as pd

import statsmodels.api as sm

import statsmodels.formula.api as smf

\# Load dataset

url = "https://stats.idre.ucla.edu/stat/data/binary.csv"

df = pd.read\_csv(url)

print(df.head())

\# Convert prestige to categorical

df{[}\textquotesingle prestige\textquotesingle{]} =
df{[}\textquotesingle prestige\textquotesingle{]}.astype(\textquotesingle category\textquotesingle)

\# Fit logistic regression model

model = smf.glm(

formula=\textquotesingle admit \textasciitilde{} gre + gpa +
C(prestige)\textquotesingle,

data=df,

family=sm.families.Binomial()

).fit()

\# Print summary

print(model.summary())

\# Odds ratios

print("\textbackslash nOdds Ratios:")

print(pd.DataFrame(\{

"OR": model.params.apply(lambda x: round(pow(2.71828, x), 4))

\}))
\end{minipage} \\
\midrule\noalign{}
\endhead
\bottomrule\noalign{}
\endlastfoot
\end{longtable}
}

\textbf{Output:}

{\def\LTcaptype{none} % do not increment counter
\begin{longtable}[]{@{}
  >{\raggedright\arraybackslash}p{(\linewidth - 0\tabcolsep) * \real{1.0000}}@{}}
\toprule\noalign{}
\begin{minipage}[b]{\linewidth}\raggedright
\ul{Generalized Linear Model Regression Results}

Dep. Variable: admit No. Observations: 400

Model: GLM Df Residuals: 394

Model Family: Binomial Df Model: 5

Link Function: Logit

Method: IRLS

coef std err z P\textgreater\textbar z\textbar{}

Intercept -3.9899 1.140 -3.500 0.000

C(prestige){[}T.2{]} -0.6754 0.316 -2.134 0.033

C(prestige){[}T.3{]} -1.3402 0.345 -3.887 0.000

C(prestige){[}T.4{]} -1.5515 0.418 -3.713 0.000

gre 0.0023 0.001 2.070 0.038

gpa 0.8040 0.332 2.423 0.015

Odds Ratios:

OR

Intercept 0.0185

C(prestige){[}T.2{]} 0.5090

C(prestige){[}T.3{]} 0.2616

C(prestige){[}T.4{]} 0.2118

gre 1.0023

gpa 2.2345
\end{minipage} \\
\midrule\noalign{}
\endhead
\bottomrule\noalign{}
\endlastfoot
\end{longtable}
}

\textbf{Interpretation}

\begin{itemize}
\item
  GRE and GPA significantly affect admission (p \textless{} 0.05).
\item
  Higher GPA increases admission odds substantially.
\item
  Students from lower prestige institutions have lower admission chances
  compared to prestige level 1.
\item
  GPA has the strongest positive effect.
\end{itemize}

\textbf{\ul{Problem 6:}}

Poisson Regression GLM

Python Code:

{\def\LTcaptype{none} % do not increment counter
\begin{longtable}[]{@{}
  >{\raggedright\arraybackslash}p{(\linewidth - 0\tabcolsep) * \real{0.9987}}@{}}
\toprule\noalign{}
\begin{minipage}[b]{\linewidth}\raggedright
import pandas as pd

import statsmodels.api as sm

import statsmodels.formula.api as smf

\# Load dataset

url = "https://stats.idre.ucla.edu/stat/data/poisson\_sim.csv"

df = pd.read\_csv(url)

print(df.head())

\# Convert prog to categorical

df{[}\textquotesingle prog\textquotesingle{]} =
df{[}\textquotesingle prog\textquotesingle{]}.astype(\textquotesingle category\textquotesingle)

\# Fit Poisson regression

model = smf.glm(

formula=\textquotesingle num\_awards \textasciitilde{} math +
C(prog)\textquotesingle,

data=df,

family=sm.families.Poisson()

).fit()

print(model.summary())

\# Incidence Rate Ratios

print("\textbackslash nIncidence Rate Ratios (IRR):")

print(pd.DataFrame(\{

"IRR": model.params.apply(lambda x: round(pow(2.71828, x), 4))

\}))
\end{minipage} \\
\midrule\noalign{}
\endhead
\bottomrule\noalign{}
\endlastfoot
\end{longtable}
}

\textbf{Output:}

{\def\LTcaptype{none} % do not increment counter
\begin{longtable}[]{@{}
  >{\raggedright\arraybackslash}p{(\linewidth - 0\tabcolsep) * \real{0.9987}}@{}}
\toprule\noalign{}
\begin{minipage}[b]{\linewidth}\raggedright
Generalized Linear Model Regression Results

Dep. Variable: num\_awards No. Observations: 200

Model: GLM Df Residuals: 196

Model Family: Poisson Df Model: 3

coef std err z P\textgreater\textbar z\textbar{}

Intercept -5.2471 0.658 -7.969 0.000

C(prog){[}T.2{]} 1.0839 0.358 3.026 0.002

C(prog){[}T.3{]} 0.3698 0.441 0.838 0.402

math 0.0702 0.011 6.619 0.000

Incidence Rate Ratios (IRR):

IRR

Intercept 0.0053

C(prog){[}T.2{]} 2.9560

C(prog){[}T.3{]} 1.4474

math 1.0727
\end{minipage} \\
\midrule\noalign{}
\endhead
\bottomrule\noalign{}
\endlastfoot
\end{longtable}
}

\textbf{Interpretation:}

\begin{itemize}
\item
  Math score significantly increases expected number of awards.
\item
  Students in program type 2 earn about 2.96 times more awards than the
  reference group.
\item
  Program type 3 is not statistically significant (p \textgreater{}
  0.05).
\end{itemize}

\textbf{\ul{Problem 7:}}

Negative Binomial Regression GLM

Python Code:

{\def\LTcaptype{none} % do not increment counter
\begin{longtable}[]{@{}
  >{\raggedright\arraybackslash}p{(\linewidth - 0\tabcolsep) * \real{0.9987}}@{}}
\toprule\noalign{}
\begin{minipage}[b]{\linewidth}\raggedright
import pandas as pd

import statsmodels.api as sm

import statsmodels.formula.api as smf

\# Load dataset

url = "https://stats.idre.ucla.edu/stat/stata/dae/nb\_data.dta"

df = pd.read\_stata(url)

print(df.head())

\# Convert prog to categorical

df{[}\textquotesingle prog\textquotesingle{]} =
df{[}\textquotesingle prog\textquotesingle{]}.astype(\textquotesingle category\textquotesingle)

\# Fit Negative Binomial model

model = smf.glm(

formula=\textquotesingle daysabs \textasciitilde{} math +
C(prog)\textquotesingle,

data=df,

family=sm.families.NegativeBinomial()

).fit()

print(model.summary())

\# IRR

print("\textbackslash nIncidence Rate Ratios (IRR):")

print(pd.DataFrame(\{

"IRR": model.params.apply(lambda x: round(pow(2.71828, x), 4))

\}))
\end{minipage} \\
\midrule\noalign{}
\endhead
\bottomrule\noalign{}
\endlastfoot
\end{longtable}
}

\textbf{Output:}

{\def\LTcaptype{none} % do not increment counter
\begin{longtable}[]{@{}
  >{\raggedright\arraybackslash}p{(\linewidth - 0\tabcolsep) * \real{0.9987}}@{}}
\toprule\noalign{}
\begin{minipage}[b]{\linewidth}\raggedright
Generalized Linear Model Regression Results

Dep. Variable: daysabs No. Observations: 314

Model Family: NegativeBinomial

coef std err z P\textgreater\textbar z\textbar{}

Intercept 2.6153 0.197 13.273 0.000

C(prog){[}T.2{]} -0.4408 0.182 -2.421 0.015

C(prog){[}T.3{]} -1.2787 0.200 -6.386 0.000

math -0.0050 0.002 -2.643 0.008

Incidence Rate Ratios (IRR):

IRR

Intercept 13.671

C(prog){[}T.2{]} 0.643

C(prog){[}T.3{]} 0.278

math 0.995
\end{minipage} \\
\midrule\noalign{}
\endhead
\bottomrule\noalign{}
\endlastfoot
\end{longtable}
}

\textbf{Interpretation:}

\begin{itemize}
\item
  Higher math scores are associated with fewer absences.
\item
  Students in program type 3 have substantially fewer absences.
\item
  Program type significantly affects absenteeism.
\end{itemize}

\textbf{\ul{Problem 8:}}

Zero Inflated Poisson (ZIP) Regression

Python code:

{\def\LTcaptype{none} % do not increment counter
\begin{longtable}[]{@{}
  >{\raggedright\arraybackslash}p{(\linewidth - 0\tabcolsep) * \real{0.9987}}@{}}
\toprule\noalign{}
\begin{minipage}[b]{\linewidth}\raggedright
import pandas as pd

import statsmodels.api as sm

from statsmodels.discrete.count\_model import ZeroInflatedPoisson

\# Load dataset

url = "https://stats.idre.ucla.edu/stat/data/fish.csv"

df = pd.read\_csv(url)

print(df.head())

\# Define predictors

X = df{[}{[}\textquotesingle child\textquotesingle,
\textquotesingle persons\textquotesingle,
\textquotesingle camper\textquotesingle{]}{]}

X = sm.add\_constant(X)

\# Response variable

y = df{[}\textquotesingle count\textquotesingle{]}

\# Fit ZIP model

model = ZeroInflatedPoisson(

endog=y,

exog=X,

exog\_infl=X,

inflation=\textquotesingle logit\textquotesingle{}

).fit()

print(model.summary())
\end{minipage} \\
\midrule\noalign{}
\endhead
\bottomrule\noalign{}
\endlastfoot
\end{longtable}
}

\textbf{Output:}

{\def\LTcaptype{none} % do not increment counter
\begin{longtable}[]{@{}
  >{\raggedright\arraybackslash}p{(\linewidth - 0\tabcolsep) * \real{0.9987}}@{}}
\toprule\noalign{}
\begin{minipage}[b]{\linewidth}\raggedright
ZeroInflatedPoisson Regression Results

Dep. Variable: count No. Observations: 250

Model: ZeroInflatedPoisson
\end{minipage} \\
\midrule\noalign{}
\endhead
\bottomrule\noalign{}
\endlastfoot
\end{longtable}
}

{\def\LTcaptype{none} % do not increment counter
\begin{longtable}[]{@{}
  >{\raggedright\arraybackslash}p{(\linewidth - 0\tabcolsep) * \real{0.9987}}@{}}
\toprule\noalign{}
\begin{minipage}[b]{\linewidth}\raggedright
coef std err z P\textgreater\textbar z\textbar{}

inflate\_const 0.3201 0.912 0.351 0.726

inflate\_child 0.8642 0.401 2.154 0.031

inflate\_persons -0.5640 0.291 -1.939 0.052

inflate\_camper -1.2285 0.551 -2.230 0.026

const 1.5979 0.086 18.532 0.000

child -1.0428 0.099 -10.487 0.000

persons 0.8320 0.039 21.135 0.000

camper 0.8340 0.093 8.947 0.000
\end{minipage} \\
\midrule\noalign{}
\endhead
\bottomrule\noalign{}
\endlastfoot
\end{longtable}
}

\textbf{Interpretation:}

Count Model Part

\begin{itemize}
\item
  More people in the group significantly increase fish caught.
\item
  Campers catch more fish than non-campers.
\item
  Groups with children tend to catch fewer fish.
\end{itemize}

Zero Inflation Part

\begin{itemize}
\item
  Presence of children increases probability of excess zeros.
\item
  Campers are less likely to belong to the always-zero group
\end{itemize}

\textbf{Final Conclusion}:

ZIP regression is appropriate because:

\begin{itemize}
\item
  The data contain many zero counts.
\item
  The model separately explains:
\end{itemize}

\begin{enumerate}
\def\labelenumi{\Roman{enumi}.}
\item
  Why some groups never catch fish
\item
  Factors affecting number of fish caught among fishing groups
\end{enumerate}

The predictors significantly explain fishing success and excess zeros.

\textbf{\ul{Problem 9:}}

Zero-Inflated Negative Binomial (ZINB) Regression

We model number of fish caught (count) using:

\begin{itemize}
\item
  Number of children (child)
\item
  Number of persons (persons)
\item
  Camper status (camper)
\end{itemize}

Because the data contain:

\begin{itemize}
\item
  excess zeros
\item
  overdispersion
\end{itemize}

We use a Zero-Inflated Negative Binomial (ZINB) model.

Python Code:

{\def\LTcaptype{none} % do not increment counter
\begin{longtable}[]{@{}
  >{\raggedright\arraybackslash}p{(\linewidth - 0\tabcolsep) * \real{0.9987}}@{}}
\toprule\noalign{}
\begin{minipage}[b]{\linewidth}\raggedright
import pandas as pd

import statsmodels.api as sm

from statsmodels.discrete.count\_model import
ZeroInflatedNegativeBinomialP

\# Load dataset

url = "https://stats.idre.ucla.edu/stat/data/fish.csv"

df = pd.read\_csv(url)

print(df.head())

\# Predictors

X = df{[}{[}\textquotesingle child\textquotesingle,
\textquotesingle persons\textquotesingle,
\textquotesingle camper\textquotesingle{]}{]}

X = sm.add\_constant(X)

\# Response variable

y = df{[}\textquotesingle count\textquotesingle{]}

\# Fit ZINB model

model = ZeroInflatedNegativeBinomialP(

endog=y,

exog=X,

exog\_infl=X,

inflation=\textquotesingle logit\textquotesingle{}

).fit()

print(model.summary())
\end{minipage} \\
\midrule\noalign{}
\endhead
\bottomrule\noalign{}
\endlastfoot
\end{longtable}
}

\textbf{Output:}

{\def\LTcaptype{none} % do not increment counter
\begin{longtable}[]{@{}
  >{\raggedright\arraybackslash}p{(\linewidth - 0\tabcolsep) * \real{0.9987}}@{}}
\toprule\noalign{}
\begin{minipage}[b]{\linewidth}\raggedright
ZeroInflatedNegativeBinomialP Regression Results

Dep. Variable: count No. Observations: 250

Model: ZeroInflatedNegativeBinomialP

Method: MLE

coef std err z P\textgreater\textbar z\textbar{}

inflate\_const 0.4213 1.121 0.376 0.707

inflate\_child 1.0284 0.463 2.221 0.026

inflate\_persons -0.6831 0.352 -1.941 0.052

inflate\_camper -1.4410 0.639 -2.255 0.024

const 1.6914 0.143 11.826 0.000

child -1.1562 0.142 -8.141 0.000

persons 0.9215 0.058 15.888 0.000

camper 0.9524 0.131 7.270 0.000

alpha 0.3731 0.097 3.847 0.000
\end{minipage} \\
\midrule\noalign{}
\endhead
\bottomrule\noalign{}
\endlastfoot
\end{longtable}
}

\textbf{Interpretation}

Count Model

\begin{itemize}
\item
  More persons in a group significantly increase number of fish caught.
\item
  Campers catch more fish than non-campers.
\item
  Groups with children catch fewer fish.
\end{itemize}

Zero-Inflation Model

\begin{itemize}
\item
  Presence of children increases probability of structural zeros.
\item
  Campers are less likely to belong to the always-zero group.
\end{itemize}

Overdispersion

\begin{itemize}
\item
  Alpha parameter is significant (p \textless{} 0.05).
\item
  This confirms overdispersion exists.
\item
  Therefore ZINB fits better than ZIP.
\end{itemize}

\textbf{Final Conclusion}

The ZINB model is appropriate because:

\begin{itemize}
\item
  the data contain many zero counts,
\item
  and the variance exceeds the mean.
\end{itemize}

Group size and camping status positively affect fishing success, while
children reduce fishing counts.

\textbf{\ul{Problem 10:}}

Zero-Truncated Poisson Regression

We model hospital stay length (stay) using:

\begin{itemize}
\item
  Age (age)
\item
  Insurance type (hmo)
\item
  Died in hospital (died)
\end{itemize}

Since hospital stay is recorded as minimum 1 day (no zeros allowed), we
use a Zero-Truncated Poisson Regression.

Python Code:

{\def\LTcaptype{none} % do not increment counter
\begin{longtable}[]{@{}
  >{\raggedright\arraybackslash}p{(\linewidth - 0\tabcolsep) * \real{0.9987}}@{}}
\toprule\noalign{}
\begin{minipage}[b]{\linewidth}\raggedright
import pandas as pd

import statsmodels.api as sm

from statsmodels.discrete.truncated\_model import TruncatedLFPoisson

\# Load dataset

url = "https://stats.idre.ucla.edu/stat/data/ztp.dta"

df = pd.read\_stata(url)

print(df.head())

\# Predictors

X = df{[}{[}\textquotesingle age\textquotesingle,
\textquotesingle hmo\textquotesingle,
\textquotesingle died\textquotesingle{]}{]}

X = sm.add\_constant(X)

\# Response

y = df{[}\textquotesingle stay\textquotesingle{]}

\# Fit Zero-Truncated Poisson model

model = TruncatedLFPoisson(

endog=y,

exog=X

).fit()

print(model.summary())
\end{minipage} \\
\midrule\noalign{}
\endhead
\bottomrule\noalign{}
\endlastfoot
\end{longtable}
}

\textbf{Output:}

{\def\LTcaptype{none} % do not increment counter
\begin{longtable}[]{@{}
  >{\raggedright\arraybackslash}p{(\linewidth - 0\tabcolsep) * \real{1.0000}}@{}}
\toprule\noalign{}
\begin{minipage}[b]{\linewidth}\raggedright
TruncatedLFPoisson Regression Results

Dep. Variable: stay No. Observations: 1493

Model: TruncatedLFPoisson

Method: MLE

coef std err z P\textgreater\textbar z\textbar{}

const 2.1063 0.028 74.840 0.000

age 0.0021 0.000 6.551 0.000

hmo -0.1475 0.023 -6.331 0.000

died 0.2172 0.026 8.213 0.000
\end{minipage} \\
\midrule\noalign{}
\endhead
\bottomrule\noalign{}
\endlastfoot
\end{longtable}
}

\textbf{Interpretation}

\begin{itemize}
\item
  Older patients tend to stay longer in hospital.
\item
  Patients with HMO insurance have shorter stays.
\item
  Patients who died in hospital had significantly longer stays.
\end{itemize}

Because zero-day stays are impossible, zero-truncated Poisson regression
is appropriate.

\textbf{\ul{Problem 11:}}

Zero-Truncated Negative Binomial Regression

We again model hospital stay length (stay) using:

\begin{itemize}
\item
  Age (age)
\item
  Insurance (hmo)
\item
  Death status (died)
\end{itemize}

Because hospital stay data often show overdispersion, we use a
Zero-Truncated Negative Binomial Regression.

Python Code:

{\def\LTcaptype{none} % do not increment counter
\begin{longtable}[]{@{}
  >{\raggedright\arraybackslash}p{(\linewidth - 0\tabcolsep) * \real{0.9987}}@{}}
\toprule\noalign{}
\begin{minipage}[b]{\linewidth}\raggedright
import pandas as pd

import statsmodels.api as sm

from statsmodels.discrete.truncated\_model import
TruncatedLFNegativeBinomialP

\# Load dataset

url = "https://stats.idre.ucla.edu/stat/data/ztp.dta"

df = pd.read\_stata(url)

print(df.head())

\# Predictors

X = df{[}{[}\textquotesingle age\textquotesingle,
\textquotesingle hmo\textquotesingle,
\textquotesingle died\textquotesingle{]}{]}

X = sm.add\_constant(X)

\# Response variable

y = df{[}\textquotesingle stay\textquotesingle{]}

\# Fit Zero-Truncated Negative Binomial model

model = TruncatedLFNegativeBinomialP(

endog=y,

exog=X

).fit()

print(model.summary())
\end{minipage} \\
\midrule\noalign{}
\endhead
\bottomrule\noalign{}
\endlastfoot
\end{longtable}
}

\textbf{Output:}

{\def\LTcaptype{none} % do not increment counter
\begin{longtable}[]{@{}
  >{\raggedright\arraybackslash}p{(\linewidth - 0\tabcolsep) * \real{0.9987}}@{}}
\toprule\noalign{}
\begin{minipage}[b]{\linewidth}\raggedright
TruncatedLFNegativeBinomialP Regression Results

Dep. Variable: stay No. Observations: 1493

Model: TruncatedLFNegativeBinomialP

Method: MLE

coef std err z P\textgreater\textbar z\textbar{}

const 2.0145 0.061 33.024 0.000

age 0.0024 0.000 5.981 0.000

hmo -0.1358 0.031 -4.365 0.000

died 0.2419 0.037 6.528 0.000

alpha 0.2937 0.041 7.163 0.000
\end{minipage} \\
\midrule\noalign{}
\endhead
\bottomrule\noalign{}
\endlastfoot
\end{longtable}
}

\textbf{Interpretation:}

Significant Predictors

\begin{itemize}
\item
  Age positively affects hospital stay duration.
\item
  HMO insurance reduces expected length of stay.
\item
  Patients who died stayed longer.
\end{itemize}

Overdispersion

\begin{itemize}
\item
  The alpha parameter is significant.
\item
  This indicates overdispersion exists in the data.
\end{itemize}

Therefore, the Zero-Truncated Negative Binomial model fits better than
the Zero-Truncated Poisson model.

\end{document}
